% =============================================================================
% capacity_measures.tex — Hospital Capacity Indicators: Methodology Reference
%
% Audience: seminar / external collaborators
% Compile:  cd Slides && TEXINPUTS=../Preambles:$TEXINPUTS xelatex -interaction=nonstopmode capacity_measures.tex
% =============================================================================

\documentclass[aspectratio=169]{beamer}
\input{header}

\usepackage{colortbl}

\title{Hospital Capacity Indicators}
\subtitle{Measurement Framework for Fondecyt 2025}
\author{Nicol\'as Libuy}
\institute{Universidad Adolfo Ib\'a\~nez}
\date{May 2026}

% =============================================================================
\begin{document}
% =============================================================================

% ── Title ────────────────────────────────────────────────────────────────────
\begin{frame}
  \titlepage
\end{frame}

% ── Motivation ───────────────────────────────────────────────────────────────
\begin{frame}{Motivation}
  \begin{columns}[T]
    \begin{column}{0.55\textwidth}
      \textbf{Research context}
      \begin{itemize}\scriptsize\setlength{\itemsep}{2pt}
        \item Fondecyt 2025 studies how severe chronic conditions (cancer, diabetes)
              affect educational and labor outcomes in Chile.
        \item Central challenge: healthcare access is endogenous to socioeconomic status
              --- patients who seek care are a selected group.
        \item Solution: exploit \emph{ex-ante} variation in hospital capacity
              \emph{before} patients are diagnosed.
      \end{itemize}
      \smallskip
      \textbf{Role of capacity indicators $z_i$}
      \begin{itemize}\scriptsize\setlength{\itemsep}{2pt}
        \item Each indicator measures the pre-diagnosis capacity environment at the
              patient's hospital.
        \item Used as a moderator in a Dynamic DiD design: high congestion $\Rightarrow$
              longer waits and deferred treatment $\Rightarrow$ worse outcomes.
      \end{itemize}
    \end{column}
    \begin{column}{0.41\textwidth}
      \begin{keybox}
        \scriptsize
        \key{Design logic:} Patients assigned to chronically saturated hospitals
        \emph{before} diagnosis face worse access regardless of their own
        characteristics. This pre-determined variation identifies treatment
        quality effects in the DiD framework.
      \end{keybox}
      \smallskip
      \begin{block}{References}
        \scriptsize
        Guidetti (2024) --- occupancy pressure \\
        Chan et al.\ (2023) --- quality benchmarks \\
        Otero (2024) --- Chilean hospitals
      \end{block}
    \end{column}
  \end{columns}
\end{frame}

% ── Data scope ───────────────────────────────────────────────────────────────
\begin{frame}{Data: MINSAL Monthly Hospital Statistics, 2014--2025}
  \begin{columns}[T]
    \begin{column}{0.50\textwidth}
      \textbf{Sources (3 datasets)}
      \begin{itemize}\scriptsize\setlength{\itemsep}{2pt}
        \item \emph{Monthly statistics} (MINSAL/DEIS 2014--2025): activity records
              per hospital $\times$ service $\times$ month.
        \item \emph{Annual bed roster / dotaci\'on} (SNSS 2010--2025):
              installed beds by hospital and functional area.
        \item \emph{Hospital registry} (DEIS 2024): location, region, complexity.
      \end{itemize}
      \smallskip
      \begin{center}\scriptsize
        \begin{tabular}{ll}
          \toprule
          Hospitals   & 205 SNSS public facilities \\
          Period      & Jan 2014 -- Nov 2025 (141 months) \\
          Services    & Codes 401--412 \\
          Group-panel & 62{,}958 hosp $\times$ group $\times$ month \\
          \bottomrule
        \end{tabular}
      \end{center}
    \end{column}
    \begin{column}{0.46\textwidth}
      \textbf{Service group definitions}
      \begin{center}\scriptsize
        \begin{tabular}{lll}
          \toprule
          Group & Codes & Description \\
          \midrule
          g1 & 401--406 & General / cancer-relevant \\
          g2 & 405--406 & High complexity (ICU) \\
          g3 & 401--412 & Broad clinical \\
          \bottomrule
        \end{tabular}
      \end{center}
      \smallskip
      \textbf{Join coverage}
      \begin{itemize}\scriptsize\setlength{\itemsep}{2pt}
        \item Dotaci\'on (capacity) matched: 95.8\,\%
        \item DEIS commune matched: 99.9\,\%
        \item Missing dotaci\'on coded \texttt{NA} --- affects capacity-normalised indicators only
      \end{itemize}
    \end{column}
  \end{columns}
\end{frame}

% ── Indicator overview ────────────────────────────────────────────────────────
\begin{frame}{Five Indicator Families}
  \begin{center}\small
    \begin{tabular}{clll}
      \toprule
      \# & Family & Variables \\
      \midrule
      1 & Capacity structure
        & \texttt{brecha\_cap},\ \texttt{bce} \\[2pt]
      2 & Occupancy \& pressure
        & \texttt{occ},\ \texttt{occ\_sat85/90/95},\ \texttt{z\_occ},\
          \texttt{sp\_3m},\ \texttt{sp\_6m} \\[2pt]
      3 & Flow \& turnover
        & \texttt{los},\ \texttt{rot},\ \texttt{egr\_x\_cama} \\[2pt]
      4 & Operational rigidity
        & \texttt{rigidez} \\[2pt]
      5 & Intensity over capacity
        & \texttt{ici} \\
      \bottomrule
    \end{tabular}
  \end{center}
  \smallskip
  \begin{keybox}
    \scriptsize
    All indicators computed at \key{Level 1} (hospital $\times$ service $\times$ month)
    and \key{Level 2} (hospital $\times$ group $\times$ month).
    Level 2 aggregates raw components by SUM, then re-derives indicators.
    Both levels exported as \texttt{.rds} + \texttt{.csv} for Stata/external linkage.
  \end{keybox}
\end{frame}

% =============================================================================
\transitionslide{1.\ Capacity Structure}
% =============================================================================

% ── brecha_cap ────────────────────────────────────────────────────────────────
\begin{frame}[shrink=7]{\texttt{brecha\_cap} --- Absolute Capacity Gap}
  \begin{columns}[T]
    \begin{column}{0.45\textwidth}
      \begin{definitionbox}{Definition}
        \[
          \text{BrechaCap}_{h,s,t}
            = D^{\text{dot}}_{h,s,a(t)} - \overline{D}_{h,s,t}
        \]
        $D^{\text{dot}}$ = installed beds (annual roster); \quad
        $\overline{D}$ = operational beds (monthly).
      \end{definitionbox}
      \smallskip
      \begin{itemize}\scriptsize\setlength{\itemsep}{3pt}
        \item \textbf{Positive:} installed capacity exceeds operations (budget,
              staffing, or maintenance constraints keep beds closed).
        \item \textbf{Negative:} temporary surge beyond the formal roster.
        \item \textbf{Units:} beds (absolute); not size-normalised ---
              use \texttt{bce} for cross-hospital comparisons.
        \item \textbf{Validation:} $\overline{D}$ correlates $r\approx 0.97$ with
              annual dotaci\'on --- they track closely but diverge in crises.
      \end{itemize}
    \end{column}
    \begin{column}{0.51\textwidth}
      \includegraphics[width=\linewidth]{../scripts/R/_outputs/fig_brecha_cap_trends.pdf}
      {\scriptsize\color{neutral}
        Group 1 monthly average, 2014--2025.
        Positive = installed $>$ operational. Dashed = zero gap.}
    \end{column}
  \end{columns}
\end{frame}

% ── bce ──────────────────────────────────────────────────────────────────────
\begin{frame}[shrink=5]{\texttt{bce} --- Relative Capacity Gap}
  \begin{columns}[T]
    \begin{column}{0.45\textwidth}
      \begin{definitionbox}{Definition}
        \[
          \text{BCE}_{h,s,t}
            = \frac{D^{\text{dot}}_{h,s,a(t)} - \overline{D}_{h,s,t}}
                   {D^{\text{dot}}_{h,s,a(t)}}
        \]
      \end{definitionbox}
      \smallskip
      \begin{itemize}\scriptsize\setlength{\itemsep}{3pt}
        \item \textbf{Interpretation:} fraction of installed beds that are
              inactive --- size-normalised for cross-hospital comparisons.
        \item \textbf{BCE $= 0.20$:} 20\,\% of installed beds not operational;
              near zero = running at full installed capacity.
        \item \textbf{Negative:} temporary expansion above the formal roster.
        \item \key{DiD role:} \normalfont size-normalised structural
              inefficiency; valid cross-hospital moderator $z_i$.
      \end{itemize}
    \end{column}
    \begin{column}{0.51\textwidth}
      \includegraphics[width=\linewidth]{../scripts/R/_outputs/fig_bce_dist.pdf}
      {\scriptsize\color{neutral}
        Distribution of BCE by period. Most hospitals near zero;
        right tail = significant inactive installed capacity.}
    \end{column}
  \end{columns}
\end{frame}

% =============================================================================
\transitionslide{2.\ Occupancy \& Pressure}
% =============================================================================

% ── occ ──────────────────────────────────────────────────────────────────────
\begin{frame}[shrink=10]{\texttt{occ} --- Occupancy Rate}
  \begin{columns}[T]
    \begin{column}{0.45\textwidth}
      \begin{definitionbox}{Definition}
        \[
          \text{Occ}_{h,s,t}
            = \frac{\text{D\'ias Cama Ocupados}_{h,s,t}}
                   {\text{D\'ias Cama Disponibles}_{h,s,t}}
        \]
      \end{definitionbox}
      \smallskip
      \begin{itemize}\scriptsize\setlength{\itemsep}{3pt}
        \item \textbf{Primary pressure measure:} fraction of operational
              bed-days occupied in month $t$.
        \item \textbf{High ($>$85\,\%):} limited buffer --- new admissions
              must displace or defer other patients.
        \item \textbf{Range:} 0.60--0.90 typical; NA when $\overline{D}\le 0$;
              0.11\,\% obs $>1$ from MINSAL rounding.
        \item \key{Primary moderator $z_i$:} \normalfont pre-diagnosis
              occupancy rate is the main candidate in the DiD specification.
      \end{itemize}
    \end{column}
    \begin{column}{0.51\textwidth}
      \includegraphics[width=\linewidth]{../scripts/R/_outputs/fig_occ_trends.pdf}
      {\scriptsize\color{neutral}
        National average by group, 2014--2025.
        Dashed = 85\,\%. Dotted = 90\,\%.
        COVID drop reflects deferred elective care.}
    \end{column}
  \end{columns}
\end{frame}

% ── occ_sat ──────────────────────────────────────────────────────────────────
\begin{frame}{\texttt{occ\_sat85/90/95} --- Binary Saturation Thresholds}
  \begin{columns}[T]
    \begin{column}{0.48\textwidth}
      \begin{definitionbox}{Definition}
        \[
          \text{OccSat}^\tau_{h,s,t}
            = \mathbf{1}[\text{Occ}_{h,s,t} > \tau]
        \]
        \centering\scriptsize$\tau \in \{0.85,\,0.90,\,0.95\}$
      \end{definitionbox}
      \smallskip
      \begin{itemize}\scriptsize\setlength{\itemsep}{3pt}
        \item \textbf{Why three thresholds:}
              85\,\% = WHO quality advisory;
              90\,\% = operational saturation (Guidetti 2024);
              95\,\% = near-total congestion.
        \item \textbf{Non-linearity:} captures step-changes in adverse
              outcomes that continuous \texttt{occ} may understate.
        \item \textbf{Input to rolling windows:} \texttt{sp\_3m}/\texttt{sp\_6m}
              are rolling means of \texttt{occ\_sat90}.
      \end{itemize}
    \end{column}
    \begin{column}{0.48\textwidth}
      \textbf{\small\% hospital-months above threshold (Group 1)}
      \smallskip
      \begin{center}\scriptsize
        \begin{tabular}{lccc}
          \toprule
          Period & $>$85\,\% & $>$90\,\% & $>$95\,\% \\
          \midrule
          Pre-COVID (14--19) & 35.2 & 22.8 & 9.2 \\
          COVID (20--22)     & 24.8 & 12.8 & 4.3 \\
          Post-COVID (23--25)& 45.4 & 27.8 & 9.8 \\
          \bottomrule
        \end{tabular}
      \end{center}
      \smallskip
      \begin{keybox}
        \scriptsize
        Post-COVID saturation \emph{exceeds} pre-COVID: pent-up demand +
        persistent workforce gaps keep beds below pre-pandemic roster.
        COVID period shows lower rates --- elective care was deferred.
      \end{keybox}
    \end{column}
  \end{columns}
\end{frame}

% ── z_occ ────────────────────────────────────────────────────────────────────
\begin{frame}{\texttt{z\_occ} --- Standardised Occupancy}
  \begin{columns}[T]
    \begin{column}{0.45\textwidth}
      \begin{definitionbox}{Definition}
        \[
          z\text{Occ}_{h,s,t}
            = \frac{\text{Occ}_{h,s,t}
                    - \mu^{\text{base}}_{h,s,m(t)}}
                   {\sigma^{\text{base}}_{h,s,m(t)}}
        \]
        $\mu^{\text{base}}, \sigma^{\text{base}}$: mean and SD of \texttt{occ}
        for the same hospital--service--\emph{calendar month}
        over 2014--2019.
      \end{definitionbox}
      \smallskip
      \begin{itemize}\scriptsize\setlength{\itemsep}{3pt}
        \item \textbf{Controls for:} hospital fixed effects and seasonal
              occupancy patterns simultaneously.
        \item \textbf{High positive:} unusually congested relative to its
              own pre-COVID historical norm.
        \item \textbf{NA} for cells with $\sigma^{\text{base}}=0$;
              empirical distribution $\approx N(0,\,0.91)$ in 2014--2019.
      \end{itemize}
    \end{column}
    \begin{column}{0.51\textwidth}
      \includegraphics[width=\linewidth]{../scripts/R/_outputs/fig_z_occ_dist.pdf}
      {\scriptsize\color{neutral}
        COVID: below-norm (left shift).
        Post-COVID: above-historical norms (right shift).}
    \end{column}
  \end{columns}
\end{frame}

% ── sp_6m ─────────────────────────────────────────────────────────────────────
\begin{frame}[shrink=10]{\texttt{sp\_3m} / \texttt{sp\_6m} --- Persistent Saturation}
  \begin{columns}[T]
    \begin{column}{0.45\textwidth}
      \begin{definitionbox}{Definition ($k \in \{3,6\}$)}
        \[
          \text{Sat}k\text{m}_{h,s,t}
            = \frac{1}{k}\sum_{j=0}^{k-1}
              \mathbf{1}[\text{Occ}_{h,s,t-j} > 0.90]
        \]
        \centering\scriptsize Trailing window; \texttt{slider::slide\_dbl}.
      \end{definitionbox}
      \smallskip
      \begin{itemize}\scriptsize\setlength{\itemsep}{3pt}
        \item \textbf{Chronic vs.\ transient:} fraction of the prior $k$ months
              above 90\,\% --- distinguishes persistent crises from spikes.
        \item \textbf{Sat6m $=1$:} above 90\,\% every month for 6 months.
        \item \textbf{Range:} [0,\,1]; mean $\approx 0.23$ for Group 1 pre-COVID.
        \item \key{Primary moderator:} \normalfont \texttt{sp\_6m} is the best
              single proxy for chronic access barriers pre-diagnosis.
      \end{itemize}
    \end{column}
    \begin{column}{0.51\textwidth}
      \includegraphics[width=\linewidth]{../scripts/R/_outputs/fig_sp6m_dist.pdf}
      {\scriptsize\color{neutral}
        Distribution of \texttt{sp\_6m} by group. Right-skewed: most
        hospitals are rarely saturated; a tail is chronically above 90\,\%.}
    \end{column}
  \end{columns}
\end{frame}

% ── Occupancy summary ─────────────────────────────────────────────────────────
\begin{frame}{Summary: Occupancy \& Pressure (Group 1)}
  \begin{center}\scriptsize
    \begin{tabular}{lcccccc}
      \toprule
      & \multicolumn{2}{c}{Pre-COVID (14--19)}
      & \multicolumn{2}{c}{COVID (20--22)}
      & \multicolumn{2}{c}{Post-COVID (23--25)} \\
      \cmidrule(lr){2-3}\cmidrule(lr){4-5}\cmidrule(lr){6-7}
      Indicator & Mean & SD & Mean & SD & Mean & SD \\
      \midrule
      \texttt{occ}     & 0.740 & 0.162 & 0.609 & 0.190 & 0.793 & 0.164 \\
      \texttt{sp\_6m}  & 0.225 & 0.319 & 0.122 & 0.273 & 0.290 & 0.353 \\
      \texttt{z\_occ}  & 0.000 & 0.912 & \muted{n/a} & \muted{n/a}
                        & \muted{n/a} & \muted{n/a} \\
      \bottomrule
    \end{tabular}
  \end{center}
  \smallskip
  \begin{keybox}
    \scriptsize
    \key{COVID pattern:} mean occ fell from 0.74 to 0.61 as elective
    procedures were cancelled and infection-control reduced operational beds.
    \key{Post-COVID rebound:} mean occ (0.79) \emph{exceeds} pre-COVID,
    driven by pent-up demand and persistent workforce shortages.
  \end{keybox}
\end{frame}

% =============================================================================
\transitionslide{3.\ Flow \& Turnover}
% =============================================================================

% ── los ──────────────────────────────────────────────────────────────────────
\begin{frame}{\texttt{los} --- Length of Stay}
  \begin{columns}[T]
    \begin{column}{0.45\textwidth}
      \begin{definitionbox}{Definition}
        \[
          \text{LOS}_{h,s,t}
            = \frac{\text{D\'ias de Estada}_{h,s,t}}
                   {\text{Egresos}_{h,s,t}}
        \]
      \end{definitionbox}
      \smallskip
      \begin{itemize}\scriptsize\setlength{\itemsep}{3pt}
        \item \textbf{What it captures:} average inpatient days per discharge
              --- determines how quickly beds become available.
        \item \textbf{High value:} slow bed turnover; delayed discharge from
              patient complexity or lack of outpatient follow-up capacity.
        \item \textbf{Cancer relevance:} long stays in upstream services
              (surgery, general medicine) delay referrals to oncology.
        \item \textbf{NA} when egresos $\le 0$; typical range 5--15 days
              (codes 401--406).
      \end{itemize}
    \end{column}
    \begin{column}{0.51\textwidth}
      \includegraphics[width=\linewidth]{../scripts/R/_outputs/fig_los_trends.pdf}
      {\scriptsize\color{neutral}
        COVID spike: selective admission of more severe cases.
        Post-COVID partial convergence toward pre-pandemic norms.}
    \end{column}
  \end{columns}
\end{frame}

% ── rot ──────────────────────────────────────────────────────────────────────
\begin{frame}{\texttt{rot} --- Bed Rotation Index}
  \begin{columns}[T]
    \begin{column}{0.45\textwidth}
      \begin{definitionbox}{Definition}
        \[
          \text{Rot}_{h,s,t}
            = \frac{\text{Egresos}_{h,s,t}}
                   {\overline{D}_{h,s,t}}
        \]
      \end{definitionbox}
      \smallskip
      \begin{itemize}\scriptsize\setlength{\itemsep}{3pt}
        \item \textbf{What it captures:} patients cycling through each
              \emph{operational} bed per month --- throughput rate.
        \item \textbf{Approximate identity:}
              $\text{Rot} \approx \text{Occ} \times D_m / \text{LOS}$,
              where $D_m$ = days in month.
        \item \textbf{High value:} rapid turnover; many patients per bed.
        \item \textbf{NA} when $\overline{D}\le 0$; typical range 1--5
              discharges/bed/month.
      \end{itemize}
    \end{column}
    \begin{column}{0.51\textwidth}
      \includegraphics[width=\linewidth]{../scripts/R/_outputs/fig_rot_trends.pdf}
      {\scriptsize\color{neutral}
        COVID collapse: lower admissions + fewer discharges (longer stays).
        Post-COVID partial recovery.}
    \end{column}
  \end{columns}
\end{frame}

% ── egr_x_cama ────────────────────────────────────────────────────────────────
\begin{frame}{\texttt{egr\_x\_cama} --- Throughput per Installed Bed}
  \begin{columns}[T]
    \begin{column}{0.45\textwidth}
      \begin{definitionbox}{Definition}
        \[
          \text{EgrXCama}_{h,s,t}
            = \frac{\text{Egresos}_{h,s,t}}
                   {D^{\text{dot}}_{h,s,a(t)}}
        \]
      \end{definitionbox}
      \smallskip
      \begin{itemize}\scriptsize\setlength{\itemsep}{3pt}
        \item \textbf{vs.\ \texttt{rot}:} rotation uses operational beds
              ($\overline{D}$); this uses installed beds ($D^{\text{dot}}$).
              The gap reflects inactive installed capacity.
        \item \textbf{High value:} high monthly throughput per formally
              allocated structural bed.
        \item \textbf{NA} when $D^{\text{dot}}\le 0$ (2.5\,\% of obs.);
              typical range 2--5 monthly discharges/installed bed.
        \item \textbf{Structural productivity:} less sensitive to short-term
              operational decisions than rotation.
      \end{itemize}
    \end{column}
    \begin{column}{0.51\textwidth}
      \includegraphics[width=\linewidth]{../scripts/R/_outputs/fig_egr_x_cama_trends.pdf}
      {\scriptsize\color{neutral}
        COVID drop: reduced admissions. Partial recovery post-2022.}
    \end{column}
  \end{columns}
\end{frame}

% ── Flow summary ──────────────────────────────────────────────────────────────
\begin{frame}{Summary: Flow \& Turnover (Group 1)}
  \begin{center}\scriptsize
    \begin{tabular}{lcccccc}
      \toprule
      & \multicolumn{2}{c}{Pre-COVID (14--19)}
      & \multicolumn{2}{c}{COVID (20--22)}
      & \multicolumn{2}{c}{Post-COVID (23--25)} \\
      \cmidrule(lr){2-3}\cmidrule(lr){4-5}\cmidrule(lr){6-7}
      Indicator & Mean & SD & Mean & SD & Mean & SD \\
      \midrule
      \texttt{los}          & 7.21 & 4.76 & 9.66 & 6.10 & 7.72 & 4.46 \\
      \texttt{rot}          & 3.45 & 1.96 & 2.10 & 1.57 & 3.60 & 1.90 \\
      \texttt{egr\_x\_cama} & 3.52 & 4.42 & 2.65 & 1.31 & 2.90 & 1.30 \\
      \bottomrule
    \end{tabular}
  \end{center}
  \smallskip
  \begin{keybox}
    \scriptsize
    \key{COVID patterns:} LOS rose sharply as elective admissions declined
    and more severe cases dominated. Rotation and throughput fell correspondingly.
    All three indicators partially recover post-2022 but remain below pre-pandemic
    levels --- reflecting persistent workforce shortages.
  \end{keybox}
\end{frame}

% =============================================================================
\transitionslide{4.\ Operational Rigidity}
% =============================================================================

% ── rigidez ──────────────────────────────────────────────────────────────────
\begin{frame}[shrink=12]{\texttt{rigidez} --- Operational Rigidity Index}
  \begin{columns}[T]
    \begin{column}{0.45\textwidth}
      \begin{definitionbox}{Definition}
        \[
          \text{Rigidez}_{h,s,t}
            = \text{Occ}_{h,s,t} \times \text{LOS}_{h,s,t}
        \]
      \end{definitionbox}
      \smallskip
      \begin{itemize}\scriptsize\setlength{\itemsep}{3pt}
        \item \textbf{Joint pressure:} combines high occupancy and long
              stays --- the two mechanisms that together block rapid bed
              turnover.
        \item \textbf{Orthogonal to occ:} corr(rigidez, occ) $\approx -0.001$
              --- adds genuinely independent information.
        \item \textbf{Cancer relevance:} surgical oncology with high LOS
              creates system-wide bottlenecks for new referrals.
        \item \key{DiD role:} \normalfont captures structural inability to
              absorb demand shocks; high-rigidez hospitals impose the
              sharpest access barriers.
      \end{itemize}
    \end{column}
    \begin{column}{0.51\textwidth}
      \includegraphics[width=\linewidth]{../scripts/R/_outputs/fig_rigidez_trends.pdf}
      {\scriptsize\color{neutral}
        COVID spike from rising LOS; higher in Group 2 (ICU).
        Post-COVID partial recovery.}
    \end{column}
  \end{columns}
\end{frame}

% =============================================================================
\transitionslide{5.\ Intensity over Capacity}
% =============================================================================

% ── ici ──────────────────────────────────────────────────────────────────────
\begin{frame}[shrink=8]{\texttt{ici} --- Intensity over Installed Capacity}
  \begin{columns}[T]
    \begin{column}{0.45\textwidth}
      \begin{definitionbox}{Definition}
        \[
          \text{ICI}_{h,s,t}
            = \frac{\text{D\'ias Cama Ocupados}_{h,s,t}}
                   {D^{\text{dot}}_{h,s,a(t)}}
        \]
      \end{definitionbox}
      \smallskip
      \begin{itemize}\scriptsize\setlength{\itemsep}{3pt}
        \item \textbf{Denominator:} installed beds ($D^{\text{dot}}$) not
              operational --- captures pressure on the full physical endowment.
        \item \textbf{Identity:} $\text{ICI} = \text{Occ}\times(1-\text{BCE})$;
              discounts occupancy by the inactive-bed fraction.
        \item \textbf{Near 1:} at or above the installed ceiling ---
              full structural utilisation.
        \item \key{vs.\ occ:} \normalfont ICI uses installed beds; occ uses
              operational. Together they triangulate both layers of the
              capacity constraint.
      \end{itemize}
    \end{column}
    \begin{column}{0.51\textwidth}
      \includegraphics[width=\linewidth]{../scripts/R/_outputs/fig_ici_trends.pdf}
      {\scriptsize\color{neutral}
        Parallels occ but lower --- systematic gap between installed and
        operational beds. COVID drop more pronounced: installed roster
        held while occupied days fell sharply.}
    \end{column}
  \end{columns}
\end{frame}

% =============================================================================
% Summary
% =============================================================================

% ── Correlation matrix ────────────────────────────────────────────────────────
\begin{frame}{Pairwise Correlations (Group 1, all periods)}
  \begin{center}\scriptsize
    \begin{tabular}{lrrrrrrrrr}
      \toprule
        & occ & sp6m & zocc & los & rot & rig. & brecha & bce & ici \\
      \midrule
      occ     & 1.000 & 0.446 & 0.219 &$-$.001& 0.379 &$-$.001&$-$.054& 0.004 & 0.262 \\
      sp6m    & 0.446 & 1.000 & 0.149 &$-$.008& 0.055 &$-$.001&$-$.121& 0.002 & 0.185 \\
      zocc    & 0.219 & 0.149 & 1.000 &$-$.004& 0.022 & 0.002 & 0.063 & 0.022 & 0.069 \\
      los     &$-$.001&$-$.008&$-$.004& 1.000 &$-$.013& 0.000 & 0.000 &$-$.018& 0.010 \\
      rot     & 0.379 & 0.055 & 0.022 &$-$.013& 1.000 &$-$.005& 0.014 & 0.026 &$-$.009\\
      rig.    &$-$.001&$-$.001& 0.002 & 0.000 &$-$.005& 1.000 & 0.002 & 0.002 &$-$.003\\
      brecha  &$-$.054&$-$.121& 0.063 & 0.000 & 0.014 & 0.002 & 1.000 & 0.394 &$-$.396\\
      bce     & 0.004 & 0.002 & 0.022 &$-$.018& 0.026 & 0.002 & 0.394 & 1.000 &$-$.933\\
      ici     & 0.262 & 0.185 & 0.069 & 0.010 &$-$.009&$-$.003&$-$.396&$-$.933& 1.000 \\
      \bottomrule
    \end{tabular}
  \end{center}
  \smallskip
  \begin{keybox}
    \scriptsize
    \key{Key findings:}
    occ--sp6m ($r=0.45$) moderate --- distinct information.
    rigidez nearly orthogonal to all occupancy measures ($|r|<0.003$) ---
    adds independent signal.
    bce--ici strongly negative ($r=-0.93$) --- algebraic complements.
  \end{keybox}
\end{frame}

% ── Data quality ──────────────────────────────────────────────────────────────
\begin{frame}{Data Quality and Scope Limitations}
  \begin{columns}[T]
    \begin{column}{0.50\textwidth}
      \textbf{Known data issues}
      \begin{itemize}\scriptsize\setlength{\itemsep}{3pt}
        \item \textbf{Occ outliers ($>$150\,\%):} 68 obs (0.11\,\%), from small rural
              facilities with $\le6$ reported available bed-days. Recommend excluding
              $\overline{D}<15$ days/month in sensitivity checks.
        \item \textbf{Missing dotaci\'on (2.5\,\%):} concentrated in 2014--2015 and
              2020--2023. Affects brecha\_cap, bce, ici, egr\_x\_cama only;
              occ, sp\_6m, los unaffected.
        \item \textbf{z-occ SD = 0.91:} cells with single pre-COVID obs
              ($\sigma=0$) are excluded from baseline, shrinking the effective pool.
      \end{itemize}
    \end{column}
    \begin{column}{0.46\textwidth}
      \textbf{Scope limitations}
      \begin{itemize}\scriptsize\setlength{\itemsep}{3pt}
        \item \textbf{Public sector only:} SNSS hospitals; private sector excluded.
        \item \textbf{Service aggregation:} Level 2 sums raw components ---
              masks service-specific variation.
        \item \textbf{COVID period:} no adjustment; recommend excluding 2020--2022
              from main DiD specifications.
        \item \textbf{Annual dotaci\'on:} intra-year changes not captured.
      \end{itemize}
      \smallskip
      \begin{keybox}
        \scriptsize
        \key{Recommendation:} use \texttt{sp\_6m} (Group 1, pre-2020 window)
        as primary moderator $z_i$; robustness with \texttt{occ} and
        \texttt{rigidez}.
      \end{keybox}
    \end{column}
  \end{columns}
\end{frame}

% =============================================================================
\end{document}
